\chapter*{Mở đầu}
\addcontentsline{toc}{chapter}{Mở đầu}

\section*{1. Giới thiệu đề tài}

Trong hoạt động của một phòng mạch tư nhân, các nghiệp vụ như tiếp nhận bệnh nhân, tra cứu hồ sơ, lập phiếu khám, kê đơn thuốc, lập hóa đơn, thanh toán và thống kê doanh thu thường xuyên diễn ra hằng ngày. Nếu các nghiệp vụ này được quản lý thủ công bằng sổ sách hoặc các tệp rời rạc, phòng mạch dễ gặp các vấn đề như trùng lặp dữ liệu bệnh nhân, khó tra cứu lịch sử khám, sai sót khi lập hóa đơn và mất thời gian khi tổng hợp báo cáo.

Từ bối cảnh đó, nhóm thực hiện đề tài \textbf{\projecttitle} nhằm xây dựng một hệ thống web nội bộ hỗ trợ quản lý các nghiệp vụ cơ bản và mở rộng của phòng mạch. Hệ thống được phát triển theo hướng client--server, trong đó frontend phục vụ giao diện người dùng, backend xử lý nghiệp vụ và cơ sở dữ liệu PostgreSQL lưu trữ thông tin tập trung.

 Phần khảo sát khách hàng được mô phỏng dựa trên nghiệp vụ phòng mạch tư nhân quy mô nhỏ đến trung bình. Hệ thống hiện hỗ trợ triển khai ở môi trường local/development, chưa có Docker, CI/CD hoặc production deployment.

\section*{2. Mục tiêu thực hiện}

Mục tiêu của đồ án là xây dựng một hệ thống có khả năng hỗ trợ các nghiệp vụ lõi của phòng mạch tư nhân và có cấu trúc đủ rõ ràng để mở rộng trong các giai đoạn tiếp theo. Các mục tiêu cụ thể gồm:

\begin{itemize}
    \item Quản lý tài khoản, đăng nhập và phân quyền người dùng theo vai trò.
    \item Quản lý hồ sơ bệnh nhân, lượt khám và lịch sử khám.
    \item Hỗ trợ bác sĩ lập phiếu khám, ghi chẩn đoán, kê đơn thuốc và hoàn tất khám.
    \item Hỗ trợ thu ngân lập hóa đơn, ghi nhận thanh toán và tra cứu hóa đơn.
    \item Hỗ trợ quản lý danh mục bệnh, thuốc, quy định và báo cáo tháng.
    \item Mở rộng quy trình vận hành với lịch hẹn, hàng đợi, sinh hiệu, dịch vụ, xét nghiệm, kho thuốc và cấp phát thuốc.
\end{itemize}

\section*{3. Nội dung thực hiện}

Nội dung thực hiện của nhóm được tổ chức theo quy trình kỹ nghệ phần mềm: khảo sát nghiệp vụ, xác định yêu cầu, thiết kế hệ thống, thiết kế phần mềm, hiện thực, kiểm thử, triển khai và vận hành. Quy trình này giúp nhóm duy trì mối liên hệ giữa yêu cầu, thiết kế, mã nguồn và kết quả kiểm thử.

\ReportFigure{figures/ch00-overview/process-overview.pdf}{Quy trình thực hiện đồ án theo định hướng kỹ nghệ phần mềm}{fig:intro-process}

\section*{4. Kế hoạch thực hiện theo hai giai đoạn}

Nhóm chia quá trình phát triển thành hai giai đoạn. Phase 1 tập trung xây dựng nghiệp vụ lõi để hệ thống có thể vận hành ở mức cơ bản. Phase 2 mở rộng hệ thống theo hướng mô phỏng quy trình vận hành phòng khám thực tế hơn.

\begin{table}[H]
\centering
\caption{Kế hoạch thực hiện theo hai giai đoạn}
\begin{tabularx}{\textwidth}{|p{2.5cm}|X|X|}
\hline
\textbf{Giai đoạn} & \textbf{Mục tiêu} & \textbf{Kết quả chính} \\
\hline
Phase 1 & Xây dựng hệ thống lõi quản lý phòng mạch tư & Auth/RBAC, bệnh nhân, lượt khám, khám bệnh, kê đơn, hóa đơn, thanh toán, danh mục, báo cáo tháng \\
\hline
Phase 2 & Mở rộng quy trình vận hành phòng khám & Lịch hẹn, hàng đợi, sinh hiệu, dịch vụ, xét nghiệm, kho thuốc, cấp phát thuốc, tổ chức phòng khám, audit log \\
\hline
\end{tabularx}
\end{table}

\section*{5. Quy trình và phương pháp quản lý nhóm}

Nhóm gồm 4 thành viên và được phân vai theo hướng chuyên trách nhưng có kiểm tra chéo. Nhóm xác nhận 4 thành viên đóng góp tương đương về công sức, dù hình thức đóng góp khác nhau theo vai trò: phân tích yêu cầu và tài liệu, thiết kế cơ sở dữ liệu và backend, thiết kế giao diện frontend, kiểm thử và tích hợp hệ thống. Git history chỉ phản ánh một phần quá trình làm việc vì nhóm có thực hiện pair programming, review trực tiếp, kiểm thử thủ công và tổng hợp tài liệu ngoài commit code. Do đó, báo cáo trình bày đóng góp dựa trên cả Git, tài liệu, test evidence, screenshot demo và biên bản phân công nhóm. Nhóm thống nhất rằng mức độ đóng góp của 4 thành viên là tương đương; sự khác biệt chủ yếu nằm ở loại công việc và bằng chứng thể hiện.

Nhóm áp dụng quy trình làm việc theo hướng \textit{Scrum-lite}: xác định backlog, chia task theo use case hoặc module, hiện thực trên branch riêng, review, tích hợp, kiểm thử và chốt baseline. Cách tổ chức này giúp nhóm kiểm soát mối liên hệ giữa yêu cầu, thiết kế, hiện thực và kiểm thử.

\section*{6. Kết quả đạt được}

Codebase hiện tại cho thấy hệ thống đạt các kết quả kỹ thuật chính sau:

\begin{table}[H]
\centering
\caption{Số liệu tổng hợp của hệ thống}
\begin{tabularx}{\textwidth}{|p{5cm}|X|}
\hline
\textbf{Hạng mục} & \textbf{Kết quả} \\
\hline
Database & 37 models, 12 enums, 3 migrations \\
\hline
Backend & 21 feature folders, 20 controllers có API route, 21 service files \\
\hline
API & 92 endpoints, gồm 41 endpoints Phase 1 và 51 endpoints Phase 2 \\
\hline
Frontend & 32 named routes, 29 page files, tổ chức theo feature \\
\hline
Testing & 3 e2e test files có ý nghĩa cho Phase 1; Phase 2 cần manual test evidence \\
\hline
Deployment & Local development only, chưa có Docker/CI/CD/production deployment \\
\hline
\end{tabularx}
\end{table}

\section*{7. Demo tổng quan hệ thống}

Phần demo tổng quan trong mở đầu nên trình bày một luồng hoàn chỉnh từ lúc bệnh nhân đến phòng mạch đến khi hoàn tất thanh toán và quản lý xem báo cáo.

\begin{table}[H]
\centering
\caption{Demo tổng quan luồng nghiệp vụ chính}
\begin{tabularx}{\textwidth}{|c|p{3.5cm}|X|p{3.5cm}|}
\hline
\textbf{Bước} & \textbf{Actor} & \textbf{Chức năng demo} & \textbf{Minh chứng cần chèn} \\
\hline
1 & Lễ tân & Tạo hồ sơ bệnh nhân và tạo lượt khám & Screenshot bệnh nhân/lượt khám \\
\hline
2 & Bác sĩ & Mở phiên khám, ghi chẩn đoán, kê đơn & Screenshot phiếu khám \\
\hline
3 & Thu ngân & Lập hóa đơn và ghi nhận thanh toán & Screenshot hóa đơn \\
\hline
4 & Quản lý & Xem báo cáo tháng & Screenshot báo cáo \\
\hline
\end{tabularx}
\end{table}

\TwoDemoFigures
{figures/screenshots/dashboard-admin.png}
{Dashboard của Admin}
{figures/screenshots/dashboard-doctor.png}
{Dashboard của Bác sĩ}
{Minh chứng giao diện tổng quan theo vai trò người dùng}
{fig:intro-dashboard-demo}
